\chapter*{СПИСОК СОКРАЩЕНИЙ И УСЛОВНЫХ ОБОЗНАЧЕНИЙ}
\addcontentsline{toc}{chapter}{СПИСОК СОКРАЩЕНИЙ И УСЛОВНЫХ ОБОЗНАЧЕНИЙ}
\setlength{\parindent}{0pt}
% Перечень располагается столбцом, слева без абзацного отступа в алфавитном порядке,
% справа через тире — расшифровка, без знаков препинания в конце строки (п. 4.8.4)
БЯМ --- большая языковая модель\\[0.5em]
CV --- Computer Vision (компьютерное зрение)\\[0.5em]
DLS --- Deep Learning School\\[0.5em]
GPU --- графический процессор\\[0.5em]
JN --- Jupyter Notebook\\[0.5em]
LMS --- Learning Management System (система управления обучением)\\[0.5em]
ML --- Machine Learning (машинное обучение)\\[0.5em]
ODS --- Open Data Science\\[0.5em]
NLP --- Natural Language Processing (обработка естественного языка)\\[0.5em]
RecSys --- recommender systems (рекомендательные системы)\\[0.5em]
SQL --- Structured Query Language (язык структурированных запросов)\\[0.5em]
ВМК --- факультет вычислительной математики и кибернетики\\[0.5em]
ДЗ --- домашнее задание\\[0.5em]
ИИ --- искусственный интеллект\\[0.5em]
КР --- контрольная работа\\[0.5em]
МО --- машинное обучение\\[0.5em]
МГУ --- Московский государственный университет\\[0.5em]
МФТИ --- Московский физико-технический институт\\[0.5em]
НИУ ВШЭ --- Национальный исследовательский университет <<Высшая школа экономики>>\\[0.5em]
ФКН --- факультет компьютерных наук\\[0.5em]
